\section{PROJETO DE INTERFACE E USABILIDADE}

\subsection{Especificação e Requisitos Estruturais das Telas}

Conforme os preceitos fundamentais da Engenharia de Software focados nas fases de especificação e projeto de produto, esta seção apresenta o levantamento conceitual detalhado e os requisitos estruturais e de navegação das cinco telas centrais do sistema \textbf{MedAgenda} \cite{shneiderman2016}. Em conformidade com a evolução de arquiteturas web interativas e com o requisito de alta usabilidade do sistema, a camada visual e o dashboard interativo foram empacotados e componentizados na pilha moderna \textbf{React.js + Vite}, garantindo altíssimo desempenho no cliente e renderização reativa sem recargas de página \cite{react2023}. O foco metodológico desta seção reside na especificação analítica e precisa dos elementos, campos transacionais e regras operacionais de interface necessários para atender aos requisitos, sem incorrer na apresentação de capturas e imagens de telas prontas \cite{nielsen1993, norman2013}.

\subsubsection{Tela 1: Autenticação de Usuários e Acesso ao Sistema}
\begin{itemize}
    \item \textbf{Objetivo do Requisito de Interface}: Viabilizar o acesso seguro, rápido e ergonômico para pacientes, médicos e administradores corporativos na plataforma, além de encaminhar novos usuários para o fluxo de cadastro (atendimento aos requisitos RF-01 e RF-02).
    \item \textbf{Elementos e Campos Obrigatórios na Estrutura}:
    \begin{itemize}
        \item Campo de entrada transacional para e-mail com validação instantânea de formato.
        \item Campo de entrada para senha com mascaramento visual de segurança e botão de visualização temporária.
        \item Botão principal de submissão para acionar a verificação de credenciais no backend.
        \item Link de direcionamento de contexto para o formulário completo de cadastro de novos pacientes.
        \item Link para o fluxo de recuperação de senha institucional ou pessoal via notificação por e-mail.
    \end{itemize}
    \item \textbf{Regras e Comportamento de Interface}: Ao autenticar com sucesso via token JWT, a aplicação SPA redireciona o usuário instantaneamente para o dashboard correspondente à sua autorização de acesso (Paciente, Médico ou Administrador). Em caso de credenciais inválidas, a interface exibe um alerta de erro explícito na parte superior sem apagar o e-mail preenchido, poupando esforço de digitação do usuário \cite{nielsen1993}.
\end{itemize}

\subsubsection{Tela 2: Agendamento de Consultas Médicas (Visão do Paciente)}
\begin{itemize}
    \item \textbf{Objetivo do Requisito de Interface}: Conduzir o paciente em um fluxo intuitivo, passo a passo e sem sobrecarga cognitiva para a marcação de consultas e acompanhamento de seu histórico clínico (atendimento aos requisitos RF-03, RF-04, RF-07 e RF-11).
    \item \textbf{Elementos e Campos Obrigatórios na Estrutura}:
    \begin{itemize}
        \item Seletor reativo em cascata listando as especialidades médicas com profissionais ativos na clínica.
        \item Seletor condicional que exibe apenas os médicos vinculados à especialidade previamente selecionada.
        \item Componente de calendário reativo para seleção da data de atendimento, bloqueando datas passadas ou feriados sem expediente.
        \item Grade dinâmica de faixas horárias livres, atualizada em tempo real com base no status da agenda do médico selecionado.
        \item Tabela transacional de agendamentos do paciente com selos visuais coloridos indicando o status atual (Confirmada, Agendada, Concluída ou Cancelada).
        \item Botão transacional de cancelamento de agendamento subordinado à verificação temporal de antecedência mínima de 24 horas.
    \end{itemize}
    \item \textbf{Regras e Comportamento de Interface}: A seleção de um horário já reservado deve ser visualmente bloqueada antes do clique para prevenir colisões \cite{norman2013}. Caso o agendamento esteja a menos de 24 horas do horário fixado, o botão de cancelamento é automaticamente desabilitado na interface, exibindo uma dica visual explicativa (\textit{Tooltip}) em conformidade estrita com o requisito RF-04.
\end{itemize}

\subsubsection{Tela 3: Agenda Diária do Médico (Visão do Médico)}
\begin{itemize}
    \item \textbf{Objetivo do Requisito de Interface}: Estruturar o fluxo de trabalho diário do profissional de saúde, organizando as chamadas cronológicas e permitindo a gestão autônoma de sua disponibilidade (atendimento aos requisitos RF-05 e RF-12).
    \item \textbf{Elementos e Campos Obrigatórios na Estrutura}:
    \begin{itemize}
        \item Componente de filtro de data rápido para navegar entre agendas passadas, do dia corrente ou agendas futuras.
        \item Filtros rápidos de segmentação por turno de atendimento clínico (Manhã ou Tarde).
        \item Tabela cronológica de pacientes contendo horário fixado, nome completo do paciente, CPF, telefone e status em tempo real do atendimento.
        \item Botão de ação primária (\textit{Atender/Prontuário}) para iniciar o atendimento e carregar o registro médico do paciente.
        \item Painel interativo para configuração e bloqueio de horários diários de disponibilidade do profissional.
    \end{itemize}
    \item \textbf{Regras e Comportamento de Interface}: A ordenação padrão dos registros é rigidamente cronológica em ordem crescente. Consultas com status Concluída permanecem listadas com destaque visual diferenciado para consulta, bloqueando qualquer edição não autorizada após a assinatura digital do prontuário \cite{shneiderman2016}.
\end{itemize}

\subsubsection{Tela 4: Registro de Prontuário Eletrônico}
\begin{itemize}
    \item \textbf{Objetivo do Requisito de Interface}: Prover uma interface ergonômica, limpa e padronizada para o registro clínico obrigatório ao término de cada atendimento médico (atendimento ao requisito RF-06).
    \item \textbf{Elementos e Campos Obrigatórios na Estrutura}:
    \begin{itemize}
        \item Cabeçalho informativo fixo exibindo dados consolidados do paciente (Nome, CPF e Idade) e dados da consulta (Data, Horário e Especialidade).
        \item Área de texto multi-linha com expansão automática para registro minucioso da anamnese, queixa principal e histórico relatado.
        \item Campo transacional padronizado para digitação do diagnóstico clínico e seleção do código na Classificação Internacional de Doenças (CID-10).
        \item Área de texto multi-linha para prescrição de medicamentos, solicitação de exames laboratoriais e orientações terapêuticas.
        \item Botão primário com confirmação de ação para assinatura digital e fechamento definitivo do prontuário eletrônico.
    \end{itemize}
    \item \textbf{Regras e Comportamento de Interface}: O preenchimento dos campos de anamnese e de diagnóstico com CID-10 é condição obrigatória para habilitar o botão de finalização. Ao assinar o documento, a aplicação altera transacionalmente o status da consulta para Concluída e redireciona o médico para a sua agenda diária \cite{nielsen1993}.
\end{itemize}

\subsubsection{Tela 5: Painel Gerencial e Indicadores (Visão da Administração)}
\begin{itemize}
    \item \textbf{Objetivo do Requisito de Interface}: Centralizar indicadores estratégicos de desempenho e viabilizar o controle administrativo sobre o corpo médico da clínica (atendimento aos requisitos RF-09 e RF-10).
    \item \textbf{Elementos e Campos Obrigatórios na Estrutura}:
    \begin{itemize}
        \item Cartões com métricas consolidadas em destaque (Total de Consultas no Período, Taxa Global de Cancelamento e Percentual de Ocupação da Clínica).
        \item Gráficos e tabelas analíticas categorizando o volume de atendimentos por especialidade médica e por status da consulta.
        \item Botões de exportação analítica permitindo baixar relatórios consolidados em formatos abertos (PDF e planilhas eletrônicas).
        \item Formulário completo de gestão para o cadastro de novos profissionais de saúde (Nome Completo, CRM, Especialidade, E-mail e Horários iniciais).
    \end{itemize}
    \item \textbf{Regras e Comportamento de Interface}: A atualização dos indicadores ocorre dinamicamente conforme o período de filtro estipulado pela administração. O cadastro de novos médicos valida previamente a duplicidade de números de CRM na base de dados \cite{nielsen1993}.
\end{itemize}

\subsection{Fluxo de Navegação e Transição de Telas}

A arquitetura e hierarquia de navegação foram concebidas com fluxos lineares, ergonômicos e livres de becos sem saída (\textit{Dead Ends}), divididos rigidamente conforme os papéis de acesso do usuário \cite{norman2013}:
\begin{itemize}
    \item \textbf{Fluxo Transacional do Paciente}: Autenticação na Plataforma ➔ Dashboard Principal de Agendamentos ➔ Seleção de Especialidade e Médico ➔ Escolha de Data e Horário Livre ➔ Confirmação e Registro do Agendamento ➔ Consulta ao Histórico Cronológico.
    \item \textbf{Fluxo Transacional do Médico}: Autenticação na Plataforma ➔ Painel da Agenda Diária Cronológica ➔ Seleção de Consulta Confirmada ➔ Preenchimento de Anamnese e CID-10 no Prontuário ➔ Assinatura Digital e Conclusão do Atendimento.
    \item \textbf{Fluxo Transacional da Administração}: Autenticação na Plataforma ➔ Dashboard Estratégico de Indicadores ➔ Gestão de Cadastros e Disponibilidade do Corpo Clínico ➔ Emissão e Exportação de Relatórios de Ocupação.
\end{itemize}

\subsection{Aplicação de Heurísticas de Usabilidade (Nielsen)}

O projeto estrutural das interfaces do MedAgenda incorpora estritamente as dez heurísticas clássicas de usabilidade de Jakob Nielsen \cite{nielsen1993, norman2013}, garantindo que a implementação final proporcione uma experiência clara, eficiente e segura para os profissionais de saúde e pacientes. A Tabela \ref{tab:nielsen} detalha essa correspondência técnica.

\begin{table}[htbp]
\centering
\small
\caption{Aplicação das Heurísticas de Nielsen na especificação estrutural da interface.}
\label{tab:nielsen}
\begin{tabularx}{\textwidth}{|>{\RaggedRight\arraybackslash}p{4cm}|>{\RaggedRight\arraybackslash}X|>{\RaggedRight\arraybackslash}p{5.2cm}|}
\hline
\textbf{Heurística de Nielsen} & \textbf{Requisito e Evidência na Especificação} & \textbf{Justificativa Técnica de Usabilidade} \\
\hline
1. Visibilidade do Status do Sistema & Exibição contínua de selos visuais diferenciados por cores para os status transacionais de consulta (Confirmada, Agendada, Concluída ou Cancelada). & Mantém o usuário plenamente ciente das mudanças de estado da sua solicitação em tempo real, reduzindo ansiedade e incertezas operacionais \cite{nielsen1993}. \\
\hline
2. Compatibilidade entre Sistema e Mundo Real & Adoção estrita da terminologia clínica e médica padronizada no Brasil (Anamnese, CID-10, Prescrição Terapêutica, CRM e Especialidade Clínica). & Respeita o modelo mental e vocabulário cotidiano de médicos, secretárias e pacientes, eliminando jargões computacionais desnecessários \cite{norman2013}. \\
\hline
3. Prevenção de Erros & Bloqueio prévio em tela de horários já reservados nos seletores e desabilitação automática da opção de cancelamento fora da regra de 24 horas. & Evita proativamente que o usuário cometa falhas contratuais ou gere colisões de horários no banco de dados \cite{shneiderman2016}. \\
\hline
4. Reconhecimento em vez de Memorização & Menus reativos em cascata listando com clareza o nome completo dos médicos disponíveis, suas especialidades clínicas e faixas horárias explícitas. & Desonera a memória de curto prazo do paciente, eliminando a necessidade de memorizar códigos internos ou consultar tabelas externas de horários \cite{nielsen1993}. \\
\hline
5. Estética Minimalista e Eficiência & Telas limpas organizadas em cartões e tabelas cronológicas, apresentando em cada visão unicamente as informações e botões acionáveis do perfil logado. & Previne a sobrecarga cognitiva e fadiga visual durante o turno diário do profissional de saúde, agilizando o atendimento clínico \cite{norman2013}. \\
\hline
\end{tabularx}\\
\autoria
\end{table}
